\begin{abstract}
Questions about scientific papers are often answered by tables and figures, yet
single-pass multimodal retrieval-augmented generation (RAG) answers every question,
supported or not, and gives the reader no way to check where a claim came from. In
this paper, an agentic multimodal RAG system is presented that gathers text, tables and
figures over several rounds, asks a sufficiency gate whether the gathered evidence is
enough before any answer is written, verifies each claim against the evidence it cites,
and draws a proof box on the page for every published claim. The system was built
with an open 4-bit 7B vision-language model on a single consumer GPU and evaluated on
150 questions from SPIQA against a single-pass and a text-only baseline. Answers were
graded by a model judge that agreed with a human grader on \nJudgeAgree{} answers
(kappa \nJudgeKappa). The agent was wrong on \nFullWrong\% of the questions, against
\nSPWrong\% for single-pass, but it was also correct less often (\nFullCorrect\% against
\nSPCorrect\%), because its abstentions did not fall on the harder questions. When the
same graph was run with a stronger model, abstention became selective. Where a claim
cited the correct table or figure, its proof box pointed at the hand-annotated answer
in \nProofInsideFull\% of cases, and every checked table-cell proof did. The gates are
therefore shown to cut confident errors, while their full benefit depends on the
capacity of the model that makes the judgements.
\end{abstract}

\begin{IEEEkeywords}
retrieval-augmented generation, multimodal question answering, agentic retrieval,
answer attribution, hallucination, scientific documents
\end{IEEEkeywords}
